\documentclass{article}
\usepackage[UTF8]{ctex}
\usepackage{amsmath,amssymb,amsfonts}
\usepackage{enumitem}
\usepackage{geometry}
\geometry{a4paper, margin=1in}
% === 错题本排版优化（不参与题目标识） ===
\setlength{\parindent}{0pt}
\setlength{\parskip}{0.35em plus 0.10em minus 0.08em}
\setlength{\emergencystretch}{3em}
\setlist{itemsep=0.25em, topsep=0.40em, parsep=0pt}
\allowdisplaybreaks[2]
% === 错题本排版优化结束 ===
\begin{document}
\section*{厦门大学2025年高等代数真题}
\section*{106．厦门大学2025年高等代数试卷}
一．填空题，每题5分，共50分。\\

\subsection{题目1}
1．\(\left(\begin{array}{cc}O & E_{3} \\ E_{2} & O\end{array}\right)^{2}=\)

\subsection{题目2}
2．设 \(A, B\) 均是可逆的4阶方阵，且 \(A\) 的每行元素之和为2，\(B\) 的每行元素之和为3，则 \(A^{-1} B+B^{-1} A\)的每行元素之和为 \(\_\_\_\_\)。

\subsection{题目3}
3．设 \(A, B\) 为 \(n\) 阶矩阵，方程组 \(A X=0, B X=0\) 分别有 \(k\) 个和 \(m\) 个线性无关的解向量，其中\(k+m>n\) ，则 \((A+B) X=0\) \(\_\_\_\_\)（有无）非零解。

\subsection{题目4}
4．设 \(V_{1}, V_{2}\) 是 \(V\) 的两个子空间，且 \(V=V_{1} \oplus V_{2}\) ，记 \(\alpha\) 是 \(V\) 中非零向量，且 \(\alpha \notin V_{1} \cap V_{2},\langle\alpha\rangle\) 是由 \(\alpha\)生成的子空间，令 \(U_{1}=V_{1} \oplus\langle\alpha\rangle, U_{2}=V_{2} \oplus\langle\alpha\rangle\) ，则 \(\operatorname{dim}\left(U_{1}+U_{2}\right)=\) \(\_\_\_\_\)。

\begin{itemize}

\subsection{题目5}
\end{itemize}

\subsection{题目5}
原题边界“5．”：设 \(\xi_{1}, \xi_{2}, \xi_{3}\) 与 \(\eta_{1}, \eta_{2}\) 分别为线性空间 \(V, W\) 的一组基，线性映射 \(\varphi: V \rightarrow W\) 满足
\[
\varphi\left(\begin{array}{lll}
\xi_{1} & \xi_{2} & \xi_{3}
\end{array}\right)=\left(\begin{array}{ll}
\eta_{1} & \eta_{2}
\end{array}\right)\left(\begin{array}{ccc}
1 & 2 & 1 \\
-1 & 1 & 0
\end{array}\right) .
\]
则 \(\operatorname{Ker} \varphi=\) \(\_\_\_\_\)．

\subsection{题目6}

\subsection{题目6}
原题边界“6．”：设 \(g(x)=x^{2}-4 x+a\) ，若存在一个唯一的首一三次多项式 \(f(x)\) ，使得 \(g(x) \mid f(x)\) ，且 \(f(x) \mid g^{2}(x)\)，则 \(f(x)=\) \(\_\_\_\_\)．

\subsection{题目7}
7．设 \(A\) 为 4 阶方阵，且 \(\operatorname{tr}\left(A^{i}\right)=4(i=1,2,3,4)\) ，则 \(|A|=\) \(\_\_\_\_\)．

\subsection{题目8}
8．\(n\) 阶方阵 \(A\) 的特征值全为2，且只有一个线性无关的特征向量，则 \(A\) 的行列式因子为 \(\_\_\_\_\)．

\subsection{题目9}
9．2024阶方阵 \(A\) 的极小多项式为 \((\lambda-1)(\lambda-2)^{2021}\) ，则 \(A\) 互不相似的 Jordan 标准形共有 \(\_\_\_\_\)种．

\subsection{题目10}
10．若实二次型 \(f\left(x_{1}, x_{2}, x_{3}\right)=a x_{1}^{2}+a x_{2}^{2}+(a-1) x_{3}^{2}+2 x_{1} x_{3}-2 x_{2} x_{3}\) 是正惯性指数为 2 的半正定型，则 \(a=\) \(\_\_\_\_\)．

二．（15 分）设 \(A^{*}=\left(\begin{array}{lll}1 & 0 & 0 \\ 0 & 9 & 0 \\ 0 & 0 & 1\end{array}\right)\) ，若 \(|A|>0\) ，且 \(A B A^{-1}=B A^{-1}+2 E\) ，求 \(B\) ．

三．（15 分）设 \(A=\left(\begin{array}{ccc}k & 1 & 1 \\ 0 & k-1 & 0 \\ 1 & 1 & k\end{array}\right), \beta=\left(\begin{array}{l}a \\ 1 \\ 1\end{array}\right)\) ，已知 \(A X=\beta\) 有无穷多解．

\begin{itemize}
  \item[（1）] 求 \(k, a\) 的值．
  \item[（2）] 求 \(A X=\beta\) 的通解．
\end{itemize}

四．（15 分）设 \(A\) 为 \(n\) 阶可逆矩阵，证明：存在 \(n-1\) 次多项式 \(g(x)\) ，使得 \(A^{-1}=g(A)\) ．

五．（15 分）设数域 \(\mathbb{F}\) 上的 \(n\) 阶矩阵 \(A\) 满足 \(|A| \neq 0\) ，且 \(A=\binom{A_{1}}{A_{2}}\) ，记 \(V_{1}, V_{2}\) 是 \(A_{1} X=0\) 和 \(A_{2} X=0\) 的解空间．证明： \(\mathbb{F}^{n}=V_{1} \oplus V_{2}\) ．

六．（15 分）设 \(A, B\) 都是正定矩阵，证明：\(A B\) 的特征值都是实数．

七．（15 分）设 \(A\in\mathbb{C}^{m\times m}, B\in\mathbb{C}^{n\times n}, C\in\mathbb{C}^{m\times n}\)，令 \(r=\operatorname{rank}(C)\)。若 \(AC=CB\)，证明：\(A,B\) 至少有 \(r\) 个相同的特征值（按代数重数计）．

八．（10 分）设 \(V\) 为复数域上 \(n\) 阶对称矩阵构成的线性空间，\(U\) 为复数域上 \(n\) 阶反称矩阵构成的线性空间，任取 \(A \in V\) ，定义

\[
\varphi: V \rightarrow U, X \mapsto A X-X A .
\]

证明：\(\varphi\) 为满射的充要条件是 \(A\) 的特征多项式在 \(\mathbb{C}\) 上无重根。
\end{document}
